(1) 充分发展动量平衡 \(0=-p_x+\partial_y\tau_{xy}\) 给出 \(\partial_y\tau_{xy}=-G\)。由中心对称性 \(\tau_{xy}(0)=0\)，故 \(\tau_{xy}=-Gy\)。若 \(Gh\leq\tau_y\)，全域应力未超过屈服值，剪切率为零；两壁无滑移使 \(u\equiv0\)。若 \(Gh>\tau_y\)，仅 \(|y|>a=\tau_y/G\) 屈服，中心 \(|y|\leq a\) 是塞状区。
(2) 对 \(a<y<h\)，因 \(u'<0\)，本构给 \(-Gy=-\tau_y+\eta u'\)，于是 \(u'=-(Gy-\tau_y)/\eta\)。从 \(u(h)=0\) 积分：
\[
u(y)=\frac{G}{2\eta}(h^2-y^2)-\frac{\tau_y}{\eta}(h-y)\quad(a\leq y\leq h).
\]
核心 \(u(y)=u(a)=G(h-a)^2/(2\eta)\)，负半区由偶对称给出。\(u'(a)=0\)，速度、剪切率连续，且两壁速度为零。
(3) 分段积分或交换积分顺序：
\[
Q=2\int_0^h u\,dy=\frac{2}{\eta}\int_a^h y(Gy-\tau_y)\,dy
=\frac{Gh^3}{3\eta}(2-3b+b^3)=\frac{Gh^3}{3\eta}(1-b)^2(2+b),\ b=a/h.
\]
\(dQ/db=(Gh^3/\eta)(b^2-1)<0\)（\(0<b<1\)）。当 \(\tau_y\to0\) 得 \(Q=2Gh^3/(3\eta)\)，正是平板Poiseuille流。
(4) 仅 \(|y|>a\) 有 \(u'\ne0\)，\(\mathcal D=2\int_a^h Gy(Gy-\tau_y)/\eta\,dy=GQ\)。也可将 \(G=-\partial_y\tau_{xy}\) 乘 \(u\) 后积分：\(\int Gu\,dy=-[\tau_{xy}u]_{-h}^h+\int\tau_{xy}u'\,dy=\mathcal D\)，边界项因无滑移为零。核心无形变故不耗散。令 \(b=1-\epsilon\)，\(Q=Gh^3\epsilon^2(3-\epsilon)/(3\eta)\sim Gh^3\epsilon^2/\eta\)，临界附近为二次开启。
